\begin{abstract}
\noindent Satellite-navigation systems and live traffic services reroute their users around congestion, so observed urban traffic increasingly reflects a mixture of drivers' intrinsic route preferences and externally-influenced routing. This matters because demand estimation, route-choice inference, congestion pricing, and infrastructure-investment decisions all rest on the assumption that observed traffic reveals what drivers actually prefer; when a significant share of trips is steered by an external recommender that itself reacts to congestion, that assumption breaks, and any downstream policy or modelling decision built on it is systematically biased. A practical tool that can distinguish externally-influenced trips from intrinsic ones, even approximately, is therefore a prerequisite for trustworthy inference on modern urban data.

This dissertation designs, validates, and stress-tests a \emph{partial-observation routing-influence filter}: a pipeline that takes a road graph, a small number of observation stations, and two snapshots of traffic taken under different conditions, fits a route-choice model to each snapshot, and scores individual trips by how likely each one is under one model versus the other. The filter is built for partial observation because urban networks contain thousands of road segments while real sensing infrastructure (loop detectors, automatic-vehicle-identification gantries, fixed cameras) typically covers only a few percent of them; any deployable tool must therefore infer behaviour over the whole network from counts at a sparse subset of locations. The inverse-Markov-chain method of Morimura, Osogami and Id\'e~\cite{morimura2013} is the route-choice estimator inside the filter, chosen because it was specifically formulated to recover a chain from such partial observations.

Existing route-choice and demand-estimation methods, from discrete-choice and recursive-logit models to inverse-reinforcement-learning approaches, are formulated for a single homogeneous behavioural regime and require either full trajectory data or strong assumptions about how to fill in unobserved segments. None of them provide a trip-level mechanism for distinguishing intrinsic routing from externally-influenced routing under sparse observation. The novel contribution of this work is to combine a partial-observation inverse-chain estimator with a trip-level log-likelihood-ratio scoring rule and a diagnostic suite that separates genuine route-choice recovery from misleading aggregate effects, producing a complete two-regime filter rather than a single-chain estimator. Three experimental settings justify and audit the design: a synthetic reproduction of the original benchmark, a controlled traffic-simulator study (SUMO~\cite{lopez2018}) in which sat-navigation rerouting is turned on and off for an identical population of vehicles, and a real-world dataset of about 10.5 million reconstructed vehicle trips from Xuancheng, China~\cite{ma2026xuancheng} using rush-hour versus off-peak travel as a behavioural stand-in for sat-navigation use. A second filter variant adds road-network attributes (road class, speed limit, lane count, turn geometry) as shared regularising priors.

The simulator study shows that the filter works end-to-end when ground-truth labels exist: the trip-level score reaches an area under the ROC curve of $0.708$ on a long-trajectory test, against a best-possible value of $0.801$ for any model of this class on the same data. The feature-regularised variant robustly improves how closely the fitted route-choice changes match the true ones at junctions, with the improvement holding across all five repeat simulations tested. A regularised two-step empirical memory test --- interpolation backoff toward the one-step reference, with count-conditional trust on the two-step counts --- lifts the empirical AUC by $+0.10$ on average across five repeat simulations (all five winning), identifying higher-order routing memory as a real design lever and motivating a parametric two-step extension to the inverter as a justified future direction; the optional hitting-rate input of the original method, by contrast, is fully tractable at a sub-network scale where its empirical estimate is informative --- where it lifts AUC by $+0.055$ at an appropriately re-tuned loss weighting and matches the full-observation reference within sampling noise --- but does not naively scale to a full city network because the wide gradient solve in the hitting-rate Jacobian benefits from dense BLAS3 kernels rather than the sparse-LU rewrite the obvious engineering fix would suggest. A simulator decomposition of the rush-hour-vs-off-peak behavioural proxy used on real-world data --- varying sat-navigation adoption and congestion level independently across five repeats --- attributes approximately $96\%$ of the peak-vs-off-peak routing-distribution signal to sat-navigation adoption rather than to congestion-induced rerouting of unaided drivers, supporting the proxy's interpretation in the controlled regime. On the real-world dataset, the filter's diagnostic suite correctly reports a negative result: once population differences between rush-hour and off-peak drivers are controlled, the filter cannot tell the two groups apart at any spatial resolution we tested, even though a full-observation reference weakly can. The contribution is therefore a designed and audited filter, together with a clear account of where it can be trusted (controlled simulation with labels) and where it currently cannot (real-world data without ground-truth labels for who is following a sat-nav).
\end{abstract}
